\wykazzrodel

% Wykaz źródeł to nie to samo co bibliografia. Bibliografia obejmuje
% opracowania naukowe, do których odsyłają powołania w tekście. Wykaz źródeł
% obejmuje narzędzia i dane, którymi się posłużyłeś: pakiety R, zbiory danych,
% serwisy udostępniające oprogramowanie.
%
% Opis pakietu R pobierz w konsoli poleceniem citation("nazwa_pakietu").

Wykaz obejmuje narzędzia programistyczne i zbiory danych wykorzystane
w projekcie.

\begin{description}
  \item[\pkg{ggplot2}] Wickham, H. (2016). \termin{ggplot2: Elegant Graphics for
    Data Analysis}. Nowy Jork: Springer.
  \item[\pkg{testthat}] Wickham, H. (2011). testthat: Get Started with Testing.
    \termin{The R Journal}, vol. 3, pp. 5-10.
  \item[\pkg{NazwaPakietu}] Nazwisko, I. (2027). \termin{NazwaPakietu: krótki
    opis przeznaczenia}. \url{https://github.com/nazwauzytkownika/NazwaPakietu}.
\end{description}
